\documentclass[11pt]{article}

\usepackage{amsmath,amssymb}
\usepackage{booktabs}
\usepackage{graphicx}
\usepackage{hyperref}
\usepackage[margin=1in]{geometry}
\usepackage{natbib}
\usepackage{lineno}

\title{Geometric Perturbation Divergence Metrics\\Are Estimation-Quality Confounded}

\author{Elliot Tower\\
\texttt{elliot@elliottower.ai}}

\date{July 2026}

\linenumbers

\begin{document}
\maketitle

\begin{abstract}
Geometric metrics for perturbation biology---direction instability,
Grassmannian geodesic distance, profile variability, optimal transport
distances---are used to quantify whether a drug or gene knockout acts
consistently across cellular contexts. We demonstrate that
variance-based, subspace-based, and distribution-based formulations of
these metrics are dominated by estimation quality: any variable that
improves the precision of the estimated object (effect magnitude,
sample size) mechanically shifts the distance metric. Across
three perturbation modalities (LINCS L1000, Perturb-seq, DepMap CRISPR
screens), we show two mechanistically distinct artifacts: (1) genes with
stronger dependency scores have mechanically higher variance across cell
lines (DepMap: raw Cohen's $d = 2.65$, coefficient-of-variation
normalized $d = -1.33$), and (2) genes with larger transcriptomic
effects produce better-estimated PCA subspaces that shift geodesic
distances regardless of directional diversity (Perturb-seq: raw $d =
0.60$, magnitude-matched $d = -0.06$, $p = 0.62$). The same
estimation-quality confound extends to optimal transport: a
pre-registered same-gene-vs-random-gene EMD test yielded $d = -0.172$,
but bin analysis revealed signal inversion across sample sizes ($d =
+0.44$ at low $n$, $d = -0.71$ at high $n$), with geodesic distance
showing even stronger inversion ($d = +1.06$ to $d = -1.24$). A
pre-registered directional prediction that essential gene knockdowns
would show greater divergence than non-essentials technically passed its
criterion ($d > 0.35$) on the naive metric; we override this
confirmation on honesty grounds after discovering the magnitude
confound. We provide a causal DAG explaining why the artifact arises,
demonstrate that magnitude-normalized metrics eliminate it, and show
that cosine-based formulations (direction instability) are inherently
immune. The corrective rule: normalize for effect magnitude before
comparing groups defined by effect strength, or geometric ``divergence''
reduces to the tautology that bigger effects are noisier.
\end{abstract}

\section{Introduction}

High-throughput perturbation screens produce multi-dimensional signatures
for thousands of drugs and gene knockdowns across panels of cell lines.
A natural question follows: does this perturbation act consistently
across contexts, or does it produce divergent effects depending on
cellular background? Geometric metrics formalize this question.
Direction instability measures the mean pairwise angular deviation of
perturbation signatures across cell lines. Grassmannian geodesic
distance measures how much a perturbation's response subspace rotates
between cell types. Profile variability (standard deviation across
contexts) measures raw spread.

These metrics are used to characterize drug mechanism conservation, gene
essentiality, and perturbation quality in LINCS L1000
\citep{subramanian2017next}, Perturb-seq \citep{replogle2022mapping},
JUMP Cell Painting \citep{chandrasekaran2023jump}, and DepMap CRISPR
screens. The implicit assumption is that ``more divergent'' reflects
biology---context-dependent mechanisms, compensatory responses, or
pathway rewiring.

We show this assumption fails for three of the four major metric
families. Variance-based metrics (standard deviation of scores
across cell lines), subspace-based metrics (geodesic distance between
PCA-estimated perturbation subspaces), and distribution-based metrics
(optimal transport / EMD) are confounded by estimation quality---any
variable that improves the precision of the estimated object (effect
magnitude, sample size) mechanically shifts the distance metric.
Cosine-based metrics (direction instability) are immune because they
normalize by magnitude before comparison.

Our evidence comes from a sequence of pre-registered and exploratory
experiments across three data modalities, reported with an explicit
integrity boundary between confirmatory and post-hoc analyses. The
discovery unifies two previously puzzling observations: (1) biological
confound corrections fail to rerank direction instability across three
domains (Spearman $\rho = 0.83$--$0.99$), and (2) essential gene
knockdowns appear to ``diverge'' more than non-essentials ($d = 0.55$ in
Perturb-seq). The first reflects magnitude-dominance of the ranking; the
second reflects magnitude-confounding of the comparison. Both dissolve
under the same DAG.

\section{Metrics and Their Magnitude Sensitivity}

\subsection{Direction instability (cosine-based, magnitude-invariant)}

For a perturbation tested in $K$ contexts with signature vectors
$\mathbf{s}_1, \ldots, \mathbf{s}_K \in \mathbb{R}^d$:
\begin{equation}
  D = 1 - \frac{2}{K(K-1)} \sum_{i < j}
    \frac{\mathbf{s}_i^\top \mathbf{s}_j}
         {\|\mathbf{s}_i\| \, \|\mathbf{s}_j\|}
\end{equation}
The $\ell_2$ normalization makes $D$ invariant to the magnitude of each
$\mathbf{s}_i$. A drug that produces a strong, consistent signal across
cell lines receives the same $D$ as one producing a weak, consistent
signal.

\subsection{Profile variability (variance-based, magnitude-sensitive)}

For a gene with dependency scores $x_1, \ldots, x_K$ across $K$ cell
lines:
\begin{equation}
  V = \text{SD}(x_1, \ldots, x_K)
\end{equation}
This metric grows with the absolute magnitude of effects. A gene with
scores uniformly near $-1$ (strongly essential everywhere) and one with
scores clustered near $0$ (non-essential everywhere) will have different
$V$ even if both are equally ``consistent''---the essential gene's
scores simply occupy a larger range.

\subsection{Grassmannian geodesic distance (subspace-based, indirectly magnitude-sensitive)}

For a gene knocked down in two cell types, let $U_1, U_2 \in
\text{Gr}(k, d)$ denote the top-$k$ PCA subspaces of the perturbation
response. The geodesic distance $\delta(U_1, U_2) =
\|\boldsymbol{\theta}\|_2$ (principal angles) is geometrically
magnitude-invariant. However, $U_1$ and $U_2$ are \emph{estimated} from
finite samples. Perturbations with larger effect sizes produce
higher-SNR data, yielding better-estimated subspaces. This estimation
quality difference propagates into the distance metric, creating an
indirect magnitude sensitivity that is absent from the population-level
definition but present in every finite-sample application.

\section{Pre-registered Experiments}

\subsection{Integrity chain}

All pre-registrations were committed before the author reviewed
experimental output (commit \texttt{249abaf}). The DepMap replication
hypothesis was committed separately (commit \texttt{98897ed}) before
execution. The analysis script was frozen (commit \texttt{92e3276})
before running. The magnitude-confound analysis (Section~\ref{sec:artifact})
is exploratory and labeled as such.

\subsection{Experiment 1: Do biological corrections rerank direction instability?}

\textbf{Hypothesis:} Removing confounding features (toxicity genes in
LINCS, essential-gene subspace in Perturb-seq, cell-health features in
JUMP-CP) would substantially reorganize the direction instability ranking
($\rho < 0.70$).

\textbf{Decision criteria:} $\rho < 0.70$: correction bites.
$0.70 \leq \rho < 0.83$: marginal. $\rho \geq 0.83$: does not bite.

\textbf{Result:} Corrections do not bite in any domain (Table~\ref{tab:robustness}).

\begin{table}[h]
\centering
\caption{Rank preservation after biological corrections across three
perturbation domains. All Spearman correlations exceed the pre-registered
``doesn't bite'' threshold of $\rho \geq 0.83$.}
\label{tab:robustness}
\begin{tabular}{llcc}
\toprule
Domain & Correction & $\rho$ & $N$ \\
\midrule
LINCS (global) & Toxicity genes removed & $>0.99$ & 8{,}949 \\
LINCS (HDAC gradient) & Toxicity genes removed & $0.83$ & 20 \\
Perturb-seq & Essential-gene subspace removed & $0.91$ & 1{,}676 \\
JUMP-CP & Cell-health features removed & $0.99$ & 25{,}254 \\
\midrule
Same-tissue LINCS & Within-tissue variance genes & $0.98$--$0.99$ & 176--358 \\
\bottomrule
\end{tabular}
\end{table}

\subsection{Experiment 2: Does essential-gene divergence replicate in DepMap?}

\textbf{Exploratory prior (disclosed):} In Perturb-seq, essential gene
knockdowns showed higher Grassmannian geodesic distance than non-essentials
(Cohen's $d = 0.55$, bootstrap CI $[0.45, 0.65]$, $p < 10^{-30}$).

\textbf{Pre-registered prediction:} Essential genes will show higher
variability of dependency profiles across cell lines than non-essential
genes (one-tailed, $d > 0.35$). Framed as confirmatory replication of
the Perturb-seq finding.

\textbf{Operationalization:} Standard deviation of PC1-corrected Chronos
scores across 1{,}178 cell lines. Essential genes defined by DepMap
common-essentials list ($n = 1{,}242$); non-essential genes defined by
DepMap nonessentials list ($n = 726$).

\textbf{Result:} $d = 2.65$ (bootstrap CI $[2.52, 2.79]$, $p <
10^{-297}$). Pre-registered criterion satisfied. PC1 removal changes
effect by 0.2\%, ruling out the fitness confound.

\section{Discovery of the Magnitude Artifact}
\label{sec:artifact}

\textbf{Status: EXPLORATORY. Not pre-registered. Discovered post-hoc
after observing that $d = 2.65$ was implausibly large.}

The $d = 2.65$ result triggered scrutiny: Cohen's $d$ values above 2
imply near-zero distributional overlap, which is unusual for biological
comparisons. Two checks revealed the effect is a magnitude artifact.

\subsection{DepMap: coefficient-of-variation normalization}

Essential genes have large negative Chronos scores (by definition---this
is what ``essential'' means in DepMap). Larger absolute scores
mechanically produce larger absolute variance. Normalizing by mean
effect magnitude (coefficient of variation: SD / $|\bar{x}|$) removes
this confound:

\begin{itemize}
  \item Essential CV: mean $= 0.38$, median $= 0.26$
  \item Non-essential CV: mean $= 2.95$, median $= 1.71$
  \item Cohen's $d$ (CV): $-1.33$ (sign reversal)
\end{itemize}

After normalization, non-essential genes show \emph{higher} relative
variability. This reversal is biologically interpretable: non-essential
genes are dispensable under most conditions but occasionally recruited
for context-specific functions, producing high relative variability
around a near-zero mean. The raw $d = 2.65$ was the tautology that
``things with bigger effects have bigger variance,'' masking a real (and
opposite-signed) pattern of context-dependent non-essential gene
recruitment.

\subsection{Perturb-seq: magnitude-matched comparison}

Geodesic distance correlates with CERES effect magnitude (Spearman $\rho
= 0.33$, $p < 10^{-40}$). After matching essential and non-essential
genes by response magnitude ($n = 127$ matched pairs within 0.2 CERES
units):

\begin{itemize}
  \item Essential mean distance: 2.659
  \item Non-essential mean distance: 2.671
  \item Cohen's $d$: $-0.06$ ($p = 0.62$, Mann-Whitney)
\end{itemize}

The effect vanishes entirely. The original $d = 0.55$ was driven by
essential genes having larger transcriptomic effects, which improves PCA
subspace estimation quality and shifts geodesic distances independent of
true directional diversity.

\subsection{Confirming the confounder with a magnitude-invariant metric}

A critical question distinguishes the confounder from the mediator
interpretation: does a mechanically magnitude-invariant metric also show
no essential/non-essential difference after magnitude control? We
computed PC1 direction instability---the cosine distance between
unit-normalized primary response directions in K562 and RPE1---for each
gene. The Perturb-seq dataset contains only two well-powered cell types
(K562 and RPE1, each with $>$50 cells per guide), so $D = 1 -
|\cos(\hat{u}_{\text{K562}}, \hat{u}_{\text{RPE1}})|$ reduces to a
single pairwise comparison rather than a population average. This
suffices for the present test: we are not estimating population-level
stability but asking whether the \emph{direction} of response differs
between these two contexts after controlling for magnitude. The metric
normalizes by construction (Equation~1), eliminating the mechanical path
from magnitude to the metric.

Raw PC1 direction instability shows a small essential/non-essential
difference ($d = 0.19$, $p = 0.001$, magnitude correlation $\rho =
0.10$). However, this residual correlation---though weaker than geodesic
distance ($\rho = 0.33$)---reflects estimation-quality effects on PCA
direction estimation. After magnitude matching ($n = 127$ pairs):

\begin{itemize}
  \item Essential PC1 D (matched): $0.505$
  \item Non-essential PC1 D (matched): $0.512$
  \item Cohen's $d$: $-0.05$ ($p = 0.63$)
\end{itemize}

After residualization (all 1{,}567 genes): $d = 0.07$, $p = 0.19$.

The magnitude-invariant metric converges to the same null as the
magnitude-sensitive one. This rules out the mediator interpretation: if
essential genes had genuinely more diverse response \emph{directions}
(independent of response strength), the cosine-based metric would detect
this regardless of magnitude matching. It does not. The confound
operates through estimation quality at every level---not just on
multi-dimensional subspaces, but on the primary PCA direction itself.

\subsection{Overriding the pre-registered confirmation}

The pre-registered criterion ($d > 0.35$) was technically satisfied by
the raw SD metric ($d = 2.65$). We override this on honesty grounds: the
metric captured magnitude, not the directional divergence it was
intended to measure. A pre-registration protects against cherry-picking
hypotheses; it does not oblige one to report a confounded metric as if
it measures the intended construct. Both the ``confirmation'' and the
artifact are reported; the latter is our actual conclusion.

\subsection{Cross-cell-type transport distances: third instance of the pattern}
\label{sec:ot_confound}

\textbf{Status: EXPLORATORY. Discovered post-hoc after running
pre-registered transport tests on Perturb-seq data (pre-registrations
frozen at commits \texttt{249abaf} and \texttt{7276811}, analysis code
SHA-verified before data examination). This is the third instance of
the same discovery pattern: a pre-registered metric yields an
apparently meaningful effect, which post-hoc bin analysis reveals as
estimation-quality confounding.}

Two pre-registered tests asked whether CRISPR knockdown responses for
the \emph{same gene} produce more similar cross-cell-type distances
than knockdowns of \emph{different genes}, testing ``gene-specific
transport'' between K562 and RPE1 \citep{replogle2022mapping}.
The first used geodesic distance on $\text{Gr}(5, 2073)$ between
top-5 PCA subspaces (1{,}676 genes); the second used earth mover's
distance (Wasserstein-1) between full perturbed-cell distributions
(1{,}875 genes, subsampled to 200 cells, 10 repeats, exact POT solver).

Both naive results showed same-gene distances lower than random-gene
distances: geodesic $d = -0.14$ ($p = 2.6 \times 10^{-6}$), EMD
$d = -0.172$ ($p = 9.8 \times 10^{-19}$). Both fell in the
pre-registered ``ambiguous'' range ($|d| < 0.2$).

Both metrics correlate strongly with minimum cell count per gene
(geodesic: Spearman $\rho = -0.84$; EMD: $\rho = -0.47$). Stratifying
by sample size reveals the same bin-inversion pattern as the DepMap
artifact (Section~\ref{sec:artifact}):

\begin{table}[h]
\centering
\caption{Cross-cell-type distance effect size inverts across
sample-size bins for both metrics. Each cell reports Cohen's $d$ of
same-gene distance relative to the random-gene null. Low-$n$ genes
appear \emph{farther} from their cross-cell-type match than random
pairs; high-$n$ genes appear closer. The aggregate effects ($d = -0.14$
geodesic, $d = -0.172$ EMD) average across bins pointing in opposite
directions. That geodesic inversion is stronger than EMD ($\rho =
-0.84$ vs.\ $-0.47$) reflects PCA subspace estimation's greater
sensitivity to sample size compared to distributional EMD estimation.
Gene counts differ between metrics (e.g., 204 vs.\ 403 in the first
bin) because the geodesic analysis used 1{,}676 genes and the EMD
analysis used 1{,}875 genes---different filtering criteria produced
different gene sets with different cell-count distributions. The
top bin ($n = 61$) has wide error bars; the monotonic trend across
all five bins is the diagnostic, and no single bin drives the
conclusion.}
\label{tab:transport_bins}
\begin{tabular}{lccrr}
\toprule
$\min(n)$ bin & $n_{\text{geo}}$ & $n_{\text{EMD}}$ & Geodesic $d$ & EMD $d$ \\
\midrule
$[20, 40)$   & 204 & 403  & $+1.06$ & $+0.44$ \\
$[40, 60)$   & 407 & 407  & $+0.58$ & $+0.02$ \\
$[60, 100)$  & 575 & 575  & $-0.14$ & $-0.27$ \\
$[100, 200)$ & 429 & 429  & $-1.13$ & $-0.71$ \\
$[200, 700)$ & 61  & 61   & $-2.02$ & $-0.75$ \\
\bottomrule
\end{tabular}
\end{table}

The monotonic inversion is the diagnostic: if gene-specific transport
were biological, all bins would show $d < 0$ regardless of cell count.
Instead, both metrics transition smoothly from positive to negative as
$n$ increases. PCA subspace estimation is more $n$-sensitive than
distributional EMD, so the metric one would naively trust more
(subspace geodesic) is more contaminated.

This extends the estimation-quality confound from
essential-vs-non-essential group comparisons
(Sections~\ref{sec:artifact}) to a same-gene-vs-random-gene transport
test. The mechanism is identical: genes with more cells yield
better-estimated objects (subspaces or distributions), producing
mechanically lower distances. The confound operates through the same
causal path drawn in Section~\ref{sec:dag}---no new mechanism is
required.

\section{Causal DAG}
\label{sec:dag}

The magnitude confound has the same causal structure across all
affected metrics, though the mechanism differs.

\subsection{DepMap (variance pathway)}

\begin{equation*}
\text{Magnitude} \rightarrow
\begin{cases}
\text{Essentiality label} & \text{(confounds grouping)} \\
\text{Score variance} & \text{(inflates metric)}
\end{cases}
\end{equation*}

Essentiality is \emph{operationally defined} by dependency-score
magnitude (Chronos $< -0.5$). The magnitude--essentiality association is
near-tautological. Magnitude also causes higher variance (larger
absolute numbers span larger ranges). Adjusting for magnitude is correct
because it removes a spurious path, not a mediating one.

\subsection{Perturb-seq (estimation pathway)}

\begin{equation*}
\text{Magnitude} \rightarrow
\begin{cases}
\text{Essentiality status} & \text{(confounds grouping)} \\
\text{Subspace estimation quality} \rightarrow \text{Geodesic distance}
  & \text{(inflates metric)}
\end{cases}
\end{equation*}

Here magnitude is not tautologically linked to essentiality, but the
confounding path through estimation quality produces the same artifact.
The evidence that this is a confounder rather than a mediator: if
magnitude were a \emph{mediator} (essential $\rightarrow$ big response
$\rightarrow$ genuinely different direction), then conditioning on
magnitude would attenuate but not eliminate the essential/non-essential
difference. Instead, magnitude-matching eliminates it completely across all tested
metrics: geodesic distance (matched $d = -0.06$), PC1 cosine D (matched
$d = -0.05$), and PC1 cosine D residualized over all genes ($d = 0.07$,
$p = 0.19$). The convergence on the same null---including the
mechanically magnitude-invariant cosine metric---confirms the entire
signal traversed the estimation-quality path.

\subsection{Cross-cell-type transport (estimation pathway, same mechanism)}

\begin{equation*}
\text{Cell count} \rightarrow
\begin{cases}
\text{Same-gene pairing structure} & \text{(confounds grouping)} \\
\text{Estimation quality} \rightarrow \text{Distance metric}
  & \text{(deflates metric)}
\end{cases}
\end{equation*}

The transport confound (Section~\ref{sec:ot_confound}) operates through
the same estimation-quality path as the Perturb-seq essential-vs-non-essential
comparison, extended to two metric families. Cell count plays the role of
``magnitude'': it determines estimation quality for both subspace-based
(geodesic, $\rho = -0.84$) and distribution-based (EMD, $\rho = -0.47$)
metrics. Same-gene pairs preserve the within-gene cell count correlation
between K562 and RPE1 ($n_{\text{K562}}$ and $n_{\text{RPE1}}$ are
positively correlated for the same gene), while random-gene pairs break
this matching. Genes with few cells produce noisy estimates that inflate
distance; genes with many cells produce converged estimates that deflate
it. The aggregate effect ($d = -0.14$ geodesic, $d = -0.172$ EMD)
averages across bins pointing in opposite directions.

\subsection{LINCS direction instability (no confound)}

\begin{equation*}
\text{Magnitude} \not\rightarrow D \quad \text{(cosine normalization
blocks the path)}
\end{equation*}

Direction instability divides by $\|\mathbf{s}_i\|$ before computing
angles. The mechanical path from magnitude to metric is blocked.
In finite samples, a residual empirical correlation persists ($\rho
\approx -0.41$ between D and signal $\ell_2$ norm in LINCS): weaker
signals produce noisier unit vectors, slightly inflating D. This
second-order effect is documented but small relative to the primary
confound in subspace metrics ($\rho = 0.33$). The cosine normalization
blocks the dominant artifact; it does not guarantee zero
magnitude-dependence in finite data.

\section{The Fix}

Three normalizations remove the magnitude confound from affected
metrics:

\begin{enumerate}
  \item \textbf{Coefficient of variation} (SD / $|\text{mean effect}|$):
    removes scale dependence. Applicable when effects have non-zero means.
  \item \textbf{Magnitude matching}: compare only perturbations with
    similar absolute effect sizes. The most assumption-free approach; does
    not require linearity.
  \item \textbf{Residualization}: regress out effect magnitude and test
    residuals. Assumes a linear (or specified) relationship between
    magnitude and the metric.
\end{enumerate}

Cosine-based formulations (direction instability, Equation~1) require no
correction---they are the magnitude-normalized version by construction.

\section{Discussion}

This paper makes one claim: variance-based, subspace-based, and
distribution-based perturbation-divergence metrics are
magnitude-confounded when used to compare groups that differ in effect
strength or sample size. The claim is supported by three instances of
the same discovery pattern: a pre-registered metric yields an apparently
meaningful effect (DepMap $d = 2.65$, Perturb-seq geodesic $d = 0.55$,
cross-cell-type transport $d = -0.14$ to $-0.172$), which post-hoc
analysis reveals as estimation-quality confounding. In each case, the
signal inverts or vanishes after controlling for the confound.

The practical implication is a rule: before comparing ``divergence''
across groups that differ in perturbation strength or cell count,
normalize for magnitude. Failing to do so produces effects that look
large ($d = 2.65$, $d = 0.55$) or directionally misleading ($d = -0.14$
averaging over bins from $+1.06$ to $-1.24$) but reduce to the
tautology that bigger effects are noisier
or better-estimated.

\subsection{Limitations}

The magnitude-matching analysis used $n = 127$ matched pairs,
constrained by the overlap in effect-size distributions between
essential and non-essential genes. These 127 pairs represent the narrow
region where both distributions overlap: atypically weak essential genes
matched to atypically strong non-essential genes. The matched sample is
not representative of either group's typical member. However, the
residualization analysis uses all 1{,}567 labeled genes and converges to
the same null ($d = 0.07$, $p = 0.19$ for PC1 D), confirming that the
matching result is not an artifact of sample restriction. Residualization
assumes a linear magnitude--D relationship; if the true relationship is
nonlinear, some confounding could survive. The magnitude-matching
approach makes no linearity assumption and yields the same null,
suggesting the linear approximation is adequate here.

Our claim that direction instability is ``immune'' to magnitude
confounding is precise: the cosine normalization removes the
\emph{mechanical} dependency. A residual empirical correlation persists
in finite samples ($\rho = 0.10$ for PC1 D vs.\ magnitude in
Perturb-seq; $\rho \approx -0.41$ between D and signal norm in LINCS
drug data). This arises because weaker signals produce noisier estimated
directions, slightly inflating D. The metric is mechanically immune but
not entirely empirically independent of magnitude when subspaces are
estimated from finite data.

Our analysis does not address whether magnitude normalization could mask
\emph{real} directional divergence in domains where essentiality is not
defined by magnitude. In principle, a biological process could produce
both larger effects AND genuinely different directions. The convergence
of four metrics (geodesic distance, subspace D, PC1 cosine D, and DepMap
SD) on the same null after magnitude control makes this unlikely in the
present data, but the possibility remains open for other systems.

\subsection{What the robustness result means in light of the artifact}

The bracket-norm robustness finding (corrections do not rerank, $\rho =
0.83$--$0.99$) is independently valid because direction instability is
cosine-based and magnitude-invariant. The robustness reflects genuine
directional stability of perturbation signatures, not magnitude
dominance. However, the \emph{interpretation} shifts: the ranking is
robust because the underlying geometry is stable, which we now
understand is partially because cosine normalization removes the primary
source of apparent variation (magnitude).

\subsection{Implications for the field}

Researchers computing variability metrics on perturbation atlases
(DepMap, Perturb-seq, PRISM, JUMP-CP) should verify that their metric is
magnitude-invariant before interpreting group differences as biological
divergence. Metrics that lack built-in normalization (SD, geodesic
distance on estimated subspaces, optimal transport distances) require
explicit magnitude correction. Metrics with built-in normalization
(cosine similarity, direction instability) are safe by construction. The
transport confound is particularly insidious because it can masquerade
as a weak positive signal: the aggregate $d = -0.14$ appeared to be
a real-but-small biological effect, when it was an average over bins
pointing in opposite directions.

\section{Conclusion}

We pre-registered divergence and transport predictions, naive metrics
confirmed them, and our own post-hoc checks killed them---three times
with the same pattern. The killing is the finding: geometric
perturbation metrics measure magnitude and sample size unless explicitly
prevented from doing so. The DAG is simple, the fix is simple, and the
failure mode is likely widespread in any analysis comparing perturbation
groups that differ in effect strength or cell count.

\section*{Data and Code Availability}

All analysis code and pre-registration documents are available at
\url{https://github.com/elliottower/bracket-norm-validity}. The
repository contains the full commit history demonstrating the
pre-registration integrity chain (commits \texttt{249abaf},
\texttt{98897ed}, \texttt{92e3276}). DepMap data (24Q4 release) is
available from \url{https://depmap.org/portal/}. LINCS L1000 data is
from GEO (GSE92742). Perturb-seq subspaces were extracted from Replogle
et al.\ (2022) supplementary data.

\section*{Author Contributions}

E.T.\ conceived the study, designed experiments, wrote pre-registrations,
implemented all analyses, discovered the magnitude artifact, and wrote
the manuscript.

\bibliographystyle{plainnat}
\bibliography{references}

\end{document}
